\begin{table}[htbp]
\centering
\caption{Electoral Cycle Effects on Crime Investigation Outcomes}
\label{tab:main}
\begin{tabular}{lcccc}
\hline\hline
 & Charge/ & No Suspect & Evid. Diff. & Log Total \\
 & Summons Rate & ID Rate & Rate & Volume \\
 & (1) & (2) & (3) & (4) \\
\hline
PCC $\times$ Pre-Election & -0.0133* & 0.0462* & 0.0036 & 0.0869 \\
 & (0.0078) & (0.0232) & (0.0044) & (0.0950) \\
PCC $\times$ Post-Election & -0.0169** & 0.0218** & 0.0023 & 0.0623 \\
 & (0.0071) & (0.0104) & (0.0074) & (0.0391) \\
\hline
Mean dep. var. & 0.107 & 0.447 & 0.324 & 9.90 \\
Force $\times$ cohort FE & Yes & Yes & Yes & Yes \\
Quarter $\times$ cohort FE & Yes & Yes & Yes & Yes \\
Observations & 1,701 & 1,701 & 1,701 & 1,701 \\
Clusters (forces) & 43 & 43 & 43 & 43 \\
\hline\hline
\end{tabular}
\begin{tablenotes}[flushleft]
\item \textit{Notes:} Stacked DiD estimates across three PCC elections (May 2016, May 2021, May 2024). Pre-Election = quarters $-4$ to $-1$ relative to election; Post-Election = quarters $0$ to $+3$. Reference period: quarter $-5$ and beyond. All specifications include force$\times$cohort and quarter$\times$cohort fixed effects. Standard errors clustered at the force level in parentheses. PCC forces ($N=41$) are compared to the Metropolitan Police and City of London Police (non-PCC controls). $^{*}p<0.10$, $^{**}p<0.05$, $^{***}p<0.01$.
\end{tablenotes}
\end{table}

